\newcommand{\DoNotLoadEpstopdf}{}
\documentclass[11pt]{article}
\usepackage[T1]{fontenc}
\usepackage{lmodern}
\usepackage{amsmath,amssymb,amsthm,mathtools}
\usepackage[margin=1.05in]{geometry}
\usepackage{microtype}
\usepackage[hidelinks]{hyperref}
\input{glyphtounicode}
\input{glyphtounicode-cmex}
\pdfgentounicode=1
\pdftrailerid{}
\hypersetup{pdftitle={A Smooth Nonquadratic Entire Affine Maximal Graph in Dimension Ten},pdfauthor={OpenAI}}
\newcommand{\R}{\mathbb R}
\DeclareMathOperator{\tr}{tr}
\newtheorem{theorem}{Theorem}[section]
\newtheorem{proposition}[theorem]{Proposition}
\newtheorem{lemma}[theorem]{Lemma}
\theoremstyle{remark}
\newtheorem{remark}[theorem]{Remark}
\title{A Smooth Nonquadratic Entire Affine Maximal Graph\\ in Dimension Ten}
\author{OpenAI}
\date{October 5, 2026}
\begin{document}
\maketitle

\begin{abstract}
We construct a smooth nonquadratic entire graph in dimension ten that
solves the classical affine maximal equation and has positive-definite
Hessian everywhere. This gives a smooth counterexample to the
entire-graph affine Bernstein assertion in dimension ten. Hessian
positivity is pointwise; no global uniform lower bound or completeness
of the Berwald--Blaschke metric is asserted.
\end{abstract}

\section{Introduction}
For a smooth function $u:\R^n\to\R$ with positive definite Hessian,
write
\begin{equation}\label{eq:affine}
 \sum_{i,j=1}^n U^{ij}w_{ij}=0,\qquad
 U^{ij}=\det(D^2u)(D^2u)^{-1}_{ij},\qquad
 w=(\det D^2u)^{-(n+1)/(n+2)}.
\end{equation}
This is the classical affine maximal equation, the Euler--Lagrange
equation of the affine area functional
$\int (\det D^2u)^{1/(n+2)}$ under compactly supported variations;
see \cite{TrudingerWang2000}.  The entire-graph form of the affine
Bernstein problem asks whether every such solution must be a
quadratic polynomial.  Throughout this paper, positivity of the
Hessian means pointwise positive definiteness, not a global lower
bound by a fixed positive matrix.

\begin{theorem}\label{thm:main}
There exists a nonquadratic function $u\in C^\infty(\R^{10})$
with $D^2u$ positive definite at every point such that
\[
 \sum_{i,j=1}^{10}U^{ij}\partial_{ij}
       (\det D^2u)^{-11/12}=0
 \quad\hbox{on }\R^{10},
 \qquad U=\det(D^2u)(D^2u)^{-1}.
\]
\end{theorem}

Thus the smooth entire-graph assertion has a negative answer in
dimension ten.  No growth assumption or completeness assumption for
the Berwald--Blaschke metric is imposed here.

The classical problem has several formulations involving different
notions of completeness.  Chern formulated the two-dimensional
entire-graph problem \cite{Chern1979}.  Calabi proved rigidity in
dimension two assuming both Euclidean completeness and completeness
of the affine metric \cite{Calabi1982}.  Trudinger and Wang later
proved that affine completeness implies Euclidean completeness for
locally uniformly convex hypersurfaces of dimension at least two
\cite{TrudingerWang2002}.
For the entire-graph problem, Trudinger and Wang proved that smooth
locally uniformly convex entire affine maximal graphs over $\R^2$
are elliptic paraboloids \cite{TrudingerWang2000}.  In the same paper,
they exhibited the singular dimension-ten example
\begin{equation}\label{eq:singular}
 u_0(x,t)=\sqrt{|x|^9+t^2},\qquad (x,t)\in\R^9\times\R,
\end{equation}
see \cite[Section~7]{TrudingerWang2000}.  This example is not smooth
at the origin and does not have positive definite Hessian everywhere.
Theorem~\ref{thm:main} supplies a smooth example satisfying precisely
these stronger pointwise requirements.  The global section-based
uniform convexity hypothesis in
\cite[Corollary~4.3]{TrudingerWang2000} is stronger than the Hessian
positivity assumed here.

Related equations with $w=(\det D^2u)^{-\theta}$ admit nonquadratic
Euclidean-complete examples for other ranges of $\theta$
\cite{Du2024,SunXingXu2026}.  The latter range extends to
$\theta=n/(n+1)$, but does not include the classical value
$\theta=11/12$ when $n=10$.

Our construction keeps the relation $u^2-t^2=F(x)^2$ suggested by
\eqref{eq:singular}, but replaces $F(x)=|x|^{9/2}$ by a smooth
strictly convex radial function with $F(0)=1$.
Section~\ref{sec:reduction} gives a dimension-independent reduction
of this ansatz to two second-order equations for positive functions
$F$ and $P$.  Section~\ref{sec:local} constructs their radial solution
smoothly through the origin.  The main obstacle is then global
existence: a priori the auxiliary function $P$ could vanish at a
finite radius.  Section~\ref{sec:global} rules this out using a bounded
invariant region for logarithmic derivatives, and completes the
proof.  The decisive change of derivative variables makes each
component of the resulting vector field nondecreasing in the other
two variables on a suitable box.  A single stationary corner then
controls all of its upper faces.

\section{Reduction to a radial system}\label{sec:reduction}
We first isolate the calculation that turns the fourth-order
equation into a system for two functions of one fewer variable.
For this section let $m\geq1$, so that the graph has $m+1$
independent variables $(x,t)\in\R^m\times\R$.  Set
\[
 k=m+2,\qquad \gamma=\frac{m+3}{m+2}.
\]
We introduce $P$ so that, for the ansatz below, the determinant power
$w$ separates into $P(x)$ times a fixed function of $t/F(x)$.

\begin{proposition}\label{prop:reduction}
Let $F,P$ be positive smooth functions on an open set
$\Omega\subset\R^m$, and suppose $H=D_x^2F$ is positive definite.
If
\begin{equation}\label{eq:reduced}
 \det H=FP^{-\gamma},\qquad
 \tr(H^{-1}D_x^2P)+kP/F=0,
\end{equation}
then
\[
 u(x,t)=\sqrt{F(x)^2+t^2}
\]
has positive definite Hessian on $\Omega\times\R$ and satisfies
\eqref{eq:affine} in dimension $n=m+1$.
\end{proposition}
\begin{proof}
Direct differentiation, interpreted as an identity of quadratic
forms, gives
\begin{equation}\label{eq:hessian}
 D^2u=\frac Fu H+\frac{F^2}{u^3}
          \left(dt-\frac tF\,dF\right)^2.
\end{equation}
The first term is positive on the $x$ directions, and the second
is positive on the remaining direction.  Thus $D^2u>0$.
Put $\zeta=t/F$ and use the pointwise basis
\[
 X_i=e_i+\zeta F_i e_t\quad(1\leq i\leq m),\qquad e_t.
\]
Its change-of-basis matrix has determinant one.  In this basis,
\eqref{eq:hessian} has diagonal blocks $(F/u)H$ and $F^2/u^3$,
so
\begin{equation}\label{eq:det}
 \det D^2u=F^{m+2}u^{-(m+3)}\det H.
\end{equation}
The first equation in \eqref{eq:reduced} is equivalent to
$P=(F/\det H)^{(m+2)/(m+3)}$.  Therefore
\begin{equation}\label{eq:w}
 w=(\det D^2u)^{-(m+2)/(m+3)}=P h(\zeta),
 \qquad h(\zeta)=(1+\zeta^2)^{k/2}.
\end{equation}

We evaluate the ordinary Hessian of $w$ in the same pointwise basis.
Since $d\zeta(X_i)=0$, twice differentiating $t=F\zeta$ gives
\[
 D^2\zeta(X_i,X_j)=-\frac\zeta F F_{ij}.
\]
Consequently
\[
 D^2w(X_i,X_j)=hP_{ij}-\frac PF\zeta h'F_{ij},
 \qquad w_{tt}=\frac P{F^2}h''.
\]
Taking the trace against the inverse blocks of $D^2u$ yields
\begin{equation}\label{eq:trace}
 \frac Fu\tr\bigl((D^2u)^{-1}D^2w\bigr)
 =h\tr(H^{-1}D_x^2P)
  +\frac PF\bigl((1+\zeta^2)h''-m\zeta h'\bigr).
\end{equation}
Differentiation of $h$ gives
\[
 (1+\zeta^2)h''-m\zeta h'
 =k(1+\zeta^2)^{k/2-1}
       \bigl(1+(k-1-m)\zeta^2\bigr)=kh,
\]
because $k=m+2$.  The second equation in \eqref{eq:reduced}
makes \eqref{eq:trace} zero.  Multiplication by the positive
determinant in \eqref{eq:det} gives \eqref{eq:affine}.
\end{proof}

We now take $F$ and $P$ radial, and use the same letters for their
profiles in $r=|x|$.  Write primes for radial derivatives and
$y=F'$.  For $r>0$ the eigenvalues of $D_x^2F$ are $y'$ in the
radial direction and $y/r$ in the $m-1$ tangential directions.
Thus, whenever $y,y'>0$,
\[
 \det D_x^2F=y'\left(\frac yr\right)^{m-1},\qquad
 \tr\bigl((D_x^2F)^{-1}D_x^2P\bigr)
 =\frac{P''}{y'}+(m-1)\frac{P'}y.
\]
Substituting into \eqref{eq:reduced}, and multiplying the second
equation by $y'y^{m-1}$, gives the equivalent system
\begin{equation}\label{eq:radial}
 y^{m-1}y'=r^{m-1}FP^{-\gamma},\qquad
 (y^{m-1}P')'=-k r^{m-1}P^{1-\gamma},\qquad F'=y.
\end{equation}
The divergence form of the second equation will allow us to start
the solution at the singular point $r=0$.

\section{A smooth radial solution at the origin}\label{sec:local}
From now on fix
\begin{equation}\label{eq:constants}
 m=9,\qquad k=11,\qquad \gamma=12/11.
\end{equation}
We construct a local solution with $F(0)=P(0)=1$.
The radial equations contain divisions by $y=F'$, which vanishes
at the origin.  An integral formulation in $s=r^2/2$ avoids this
singularity and makes smoothness in $x$ explicit.

\begin{lemma}\label{lem:local}
There is $\varepsilon>0$ and a solution $F,P$ of
\eqref{eq:radial} on $0<r<\varepsilon$ such that
$F(|x|)$ and $P(|x|)$ extend smoothly to $|x|<\varepsilon$,
with
\[
 F(0)=P(0)=1,\qquad D_x^2F(0)=I,\qquad
 D_x^2P(0)=-\frac{11}{9}I.
\]
Moreover $F,P,y,y'>0$ and $P'<0$ for $0<r<\varepsilon$.
\end{lemma}
\begin{proof}
Seek $F(r)=f(s)$ and $P(r)=p(s)$, where $s=r^2/2$.
The integrated equations suggest the normalized derivative
$q=y/r$ and flux $j=-y^{m-1}P'/(kr^m)$ for $r>0$.
We will construct their regular extensions as functions of $s$
using the following averages.
For continuous $f,p$ on $[-\delta,\delta]$ taking values in
$[1/2,3/2]$, define
\begin{align}
 q(s)&=\left(m\int_0^1 v^{m-1}
                  f(v^2s)p(v^2s)^{-\gamma}\,dv\right)^{1/m},
                  \label{eq:q}\\
 j(s)&=\int_0^1 v^{m-1}p(v^2s)^{1-\gamma}\,dv.
                  \label{eq:j}
\end{align}
On the closed subset
\[
 \mathcal K=\{(f,p)\in C([-\delta,\delta])^2:
          \|f-1\|_\infty,\|p-1\|_\infty\leq1/2\},
\]
consider
\begin{equation}\label{eq:fixed-point}
 \mathcal T(f,p)(s)=
 \left(1+\int_0^s q(a)\,da,
       \ 1-k\int_0^s j(a)q(a)^{1-m}\,da\right).
\end{equation}
The functions $q,j$ have uniform positive upper and lower bounds.
The maps $(f,p)\mapsto q$ and $(f,p)\mapsto jq^{1-m}$ are
Lipschitz in the product supremum norm, with constants independent
of $\delta$: all scalar powers are evaluated on fixed positive
compact intervals.  Hence there are constants $M,L$ independent
of $\delta$ such that
\[
 \|\mathcal T(f,p)-(1,1)\|_\infty\leq M\delta,
 \qquad
 \|\mathcal T(f,p)-\mathcal T(\widetilde f,\widetilde p)\|_\infty
 \leq L\delta\|(f,p)-(\widetilde f,\widetilde p)\|_\infty.
\]
Choose $M\delta\leq1/2$ and $L\delta<1$.
The contraction mapping theorem gives a fixed point in $\mathcal K$.

The fixed point is smooth on $(-\delta,\delta)$.
Indeed, continuity of $q,j$ first gives $f,p\in C^1$.
If $f,p$ are $C^\ell$, differentiation under the integrals in
\eqref{eq:q}--\eqref{eq:j} shows that $q,j$ are $C^\ell$,
and \eqref{eq:fixed-point} gives one further derivative.
This proves the assertion by induction.  The definition on both
sides of $s=0$ ensures that $f(|x|^2/2)$ and $p(|x|^2/2)$
are smooth in $x$ at the origin.

For $r>0$, the fixed point equations give $y=rq(r^2/2)>0$.
Changing variables $\rho=rv$ in \eqref{eq:q}--\eqref{eq:j}
then yields
\begin{equation}\label{eq:integral}
 y^m=m\int_0^r \rho^{m-1}F(\rho)P(\rho)^{-\gamma}\,d\rho,
 \qquad
 y^{m-1}P'=-k\int_0^r\rho^{m-1}P(\rho)^{1-\gamma}\,d\rho.
\end{equation}
Differentiation proves \eqref{eq:radial}; the first equation
implies $y'>0$, and the second integral gives $P'<0$.
Finally $q(0)=1$ and $j(0)=1/m$, so
\[
 f_s(0)=1,\qquad p_s(0)=-k/m=-11/9.
\]
These are exactly the claimed Hessians of the radial extensions.
Shrinking $\varepsilon$ if necessary completes the proof.
\end{proof}

\section{An invariant box and global continuation}\label{sec:global}
The local solution now has every required property.  It remains
to show that it can be continued to arbitrarily large radii without
losing positivity.  We do this by controlling logarithmic derivatives,
rather than by estimating $F$ and $P$ separately.

Take the maximal outward continuation of the local solution on
which $F,P,y>0$.  This is a regular smooth ODE for $r>0$.
Explicitly, with $Z=P'$ and
\[
 G(r,F,P,y)=\frac{r^8FP^{-12/11}}{y^8},
\]
it reads
\begin{equation}\label{eq:regular}
 F'=y,\qquad P'=Z,\qquad y'=G,\qquad
 Z'=-8\frac Gy Z-\frac{11r^8P^{-1/11}}{y^8}.
\end{equation}
In particular $y'>0$ throughout this continuation.
Define
\begin{equation}\label{eq:ratios}
 A=\frac{ry}{F},\qquad B=\frac{ry'}y,
 \qquad C=-\frac{rP'}P,
\end{equation}
and let a dot mean differentiation with respect to $\log r$.
Since $B>0$, we may also set $D=C/B$.
This normalization gives a vector field with nonnegative cross
derivatives on the box in Lemma~\ref{lem:box}, allowing its upper
faces to be controlled by one corner.

\begin{lemma}\label{lem:autonomous}
On the interval of continuation, $A,B>0$, and the functions
$A,B,D$ satisfy
\begin{equation}\label{eq:autonomous}
 \begin{aligned}
  \dot A&=A(1+B-A),\\
  \dot B&=B\bigl(9+A+(12/11)BD-9B\bigr),\\
  \dot D&=(11-D)A-8D+BD(1-D/11).
 \end{aligned}
\end{equation}
Moreover $(A,B,D)\to(0,1,0)$ as $r\downarrow0$, and all three
coordinates are positive for sufficiently small $r>0$.
\end{lemma}
\begin{proof}
Positivity of $F,y,y'$ gives $A,B>0$.
Using \eqref{eq:radial}, one may write
\[
 B=\frac{r^mFP^{-\gamma}}{y^m}.
\]
Logarithmic differentiation of $A$ and $B$, followed by
the second equation in \eqref{eq:radial}, gives
\begin{align*}
 \dot A&=A(1+B-A),\\
 \dot B&=B(m+A+\gamma C-mB),\\
 \dot C&=C(1+C-(m-1)B)+kAB.
\end{align*}
For the last identity, use
$\dot C=C+C^2-r^2P''/P$ and
$AB=r^{m+1}P^{-\gamma}/y^{m-1}$.
Thus
\[
 \dot D=(k-D)A-(m-1)D+BD\bigl(1+(1-\gamma)D\bigr).
\]
Substituting \eqref{eq:constants} proves \eqref{eq:autonomous}.

Lemma~\ref{lem:local} gives
$F=1+O(r^2)$, $y=r+O(r^3)$, $y'=1+O(r^2)$,
and $P'=-(11/9)r+O(r^3)$.
Hence $A\to0$, $B\to1$, and $D\to0$.
The positivity of $A,B$ and the local inequality $P'<0$
also give $D>0$ for small positive $r$.
\end{proof}

The singular profile \eqref{eq:singular} suggests the bounds for
these variables.  Its radial factor $F=r^{9/2}$ has
$A=9/2$ and $B=7/2$; the associated function
$P=(F/\det D_x^2F)^{11/12}$ is a positive constant times
$r^{-33/2}$, giving $C=33/2$ and $D=33/7$.
These constants determine an invariant box for the smooth solution.

\begin{lemma}\label{lem:box}
The box
\begin{equation}\label{eq:box}
 \mathcal B=[0,9/2]\times[0,7/2]\times[0,33/7]
\end{equation}
is forward invariant for \eqref{eq:autonomous}.
The solution arising from Lemma~\ref{lem:local} lies in
$\mathcal B$ throughout its interval of continuation.
\end{lemma}
\begin{proof}
Write $V=(V_1,V_2,V_3)$ for the vector field in
\eqref{eq:autonomous}.  The upper corner
$a_*=(9/2,7/2,33/7)$ satisfies $V(a_*)=0$ by direct substitution.
Each $V_i$ is nondecreasing in the other two coordinates on
$\mathcal B$.  The nonzero cross derivatives are
\[
 \partial_BV_1=A,\quad \partial_AV_2=B,\quad
 \partial_DV_2=(12/11)B^2,\quad
 \partial_AV_3=11-D,\quad
 \partial_BV_3=D(1-D/11),
\]
all nonnegative because $0\leq D\leq33/7<11$.
On each upper face, comparison with $a_*$ therefore gives
the corresponding $V_i\leq0$.
On the lower faces $A=0$, $B=0$, and $D=0$, the respective
components are $0$, $0$, and $11A$, hence nonnegative.

These weak boundary inequalities suffice for invariance.
Let $\Pi z$ be the coordinatewise projection of
$z\in\R^3$ onto $\mathcal B$.  The face inequalities imply
\[
 (z-\Pi z)\cdot V(\Pi z)\leq0.
\]
On any bounded neighborhood of the box, $V$ has a Lipschitz
constant $L$.  Along a trajectory $z$ in that neighborhood,
the squared distance $d^2=|z-\Pi z|^2$ consequently satisfies
\[
 \frac{d}{d\log r}d^2
 =2(z-\Pi z)\cdot V(z)\leq2L d^2.
\]
The squared distance to a closed box is continuously differentiable,
so this identity also holds when the nearest face changes.
Gronwall's inequality shows that a trajectory starting in the box
cannot leave it.

Finally, Lemma~\ref{lem:autonomous} gives strictly positive
coordinates tending to $(0,1,0)$ as $r\downarrow0$.
They therefore lie below all three upper bounds for sufficiently
small $r>0$.  Forward invariance proves the last assertion.
\end{proof}

\begin{proposition}\label{prop:global}
The solution in Lemma~\ref{lem:local} extends to every $r>0$,
with $F,P,y,y'>0$.  Its radial extensions to $\R^9$ are smooth,
and $D_x^2F$ is positive definite everywhere.
\end{proposition}
\begin{proof}
Suppose its maximal outer endpoint were $R<\infty$, and fix
$r_0\in(0,R)$.  Lemma~\ref{lem:box} gives
\[
 0\leq A\leq9/2,\qquad 0\leq B\leq7/2,
 \qquad 0\leq C=BD\leq33/2.
\]
Since the derivatives of $\log F,\log y,\log P$ with respect
to $\log r$ are $A,B,-C$, respectively, integration gives,
for $r_0\leq r<R$,
\begin{align*}
 F(r_0)&\leq F(r)\leq F(r_0)(R/r_0)^{9/2},\\
 y(r_0)&\leq y(r)\leq y(r_0)(R/r_0)^{7/2},\\
 P(r_0)(r_0/R)^{33/2}&\leq P(r)\leq P(r_0).
\end{align*}
Also $|P'(r)|=C(r)P(r)/r\leq(33/2)P(r_0)/r_0$.
Thus $(r,F,P,y,Z)$ remains in a compact subset of the domain
$r,F,P,y>0$ of the smooth system \eqref{eq:regular}.
Its right-hand side is bounded there, so the state has a limit
at $R$ in that domain.  Local existence and uniqueness at this
limit extend the solution past $R$, a contradiction.

Smoothness for $r>0$ follows from \eqref{eq:regular}, and
Lemma~\ref{lem:local} supplies smoothness in $x$ at the origin.
The Hessian eigenvalues for $r>0$ are $y'>0$ and $y/r>0$;
at the origin the Hessian is $I$.
\end{proof}

\begin{proof}[Proof of Theorem~\ref{thm:main}]
Take the global radial functions of Proposition~\ref{prop:global}.
By \eqref{eq:radial}, they satisfy \eqref{eq:reduced} for $x\ne0$.
The functions are smooth, $F,P>0$, and $D_x^2F>0$ everywhere,
so both identities extend to $x=0$ by continuity.
Proposition~\ref{prop:reduction} now shows that
\[
 u(x,t)=\sqrt{F(|x|)^2+t^2},\qquad (x,t)\in\R^9\times\R,
\]
is smooth, has positive definite Hessian everywhere, and solves
the classical affine maximal equation in dimension ten.
Finally, $u(0,t)=\sqrt{1+t^2}$ because $F(0)=1$,
so $u$ is not a quadratic polynomial.
\end{proof}

{\small
\bibliographystyle{abbrv}
\bibliography{references}
}
\end{document}
